\documentclass[12pt]{article}
\usepackage[T1]{fontenc}
\usepackage[utf8]{inputenc}
\usepackage{lmodern}
\usepackage{textcomp}
\usepackage{enumitem}
\usepackage{geometry}
\geometry{margin=1in}
\begin{document}
\begin{center}
Answer Key - Political Science - Graduate Level Assessment 10 \
\small (Answers are ordered exactly as in the question paper.)
\end{center}

\section*{Multiple Choice}
\begin{enumerate}[label=\arabic*.]
\item \textbf{Answer:} A
\\ \textit{Options:} A) Provide strong enforcement; B) Monitor parties; C) Provide fora for discussion; D) Reduce transaction costs for agreements
\\ \textit{Note:} Intergovernmental organizations seldom provide strong enforcement because they typically rely on the voluntary cooperation of member states and lack the authority or resources to enforce compliance through coercive means.
\item \textbf{Answer:} A
\\ \textit{Options:} A) George Kennan; B) John Foster Dulles; C) Henry Kissinger; D) Dwight Eisenhower
\\ \textit{Note:} George Kennan is considered the "father" of containment because he articulated the strategy in his 1946 "Long Telegram," which outlined the need for the United States to counter Soviet expansionism through a combination of military, economic, and diplomatic measures.
\item \textbf{Answer:} B
\\ \textit{Options:} A) Jimmy Carter; B) Ronald Reagan; C) Bill Clinton; D) Dwight Eisenhower
\\ \textit{Note:} The correct answer is Ronald Reagan because his administration significantly escalated military spending during the 1980 s as part of a broader strategy to combat the Soviet Union and enhance U.S. defense capabilities, leading to the largest peacetime defense budget in American history.
\item \textbf{Answer:} C
\\ \textit{Options:} A) the importance of military superiority.; B) how the importance of oil is often overexaggerated.; C) the increasing importance of economic instruments of foreign policy.; D) the need to drill for domestic sources.
\\ \textit{Note:} The correct answer emphasizes that oil's strategic significance has led the U.S. to prioritize economic tools, such as sanctions and trade agreements, in its foreign policy to secure energy resources and maintain global economic dominance.
\item \textbf{Answer:} A
\\ \textit{Options:} A) Its general anti-statism opposes centralized government; B) It views world government as impractical; C) It views the UN as a sufficient form of global governance; D) None of the above
\\ \textit{Note:} American exceptionalism, which emphasizes individualism and limited government, inherently promotes a skepticism towards centralized authority, thereby fostering opposition to world government structures that may be perceived as infringing on national sovereignty and personal freedoms.
\end{enumerate}

\section*{True / False}
\begin{enumerate}[label=\arabic*.]
\setcounter{enumi}{5}
\item \textbf{Answer:} True
\\ \textit{Options:} A) True; B) False
\\ \textit{Note:} The Marshall Plan provided substantial financial assistance to European nations after World War II to promote economic stability, thereby reducing the appeal of communism and countering Soviet influence in the region.
\item \textbf{Answer:} False
\\ \textit{Options:} A) True; B) False
\\ \textit{Note:} Democratic enlargement primarily aimed at promoting democracy in authoritarian regimes and expanding democratic governance globally, rather than restructuring the internal political systems of existing democracies.
\item \textbf{Answer:} False
\\ \textit{Options:} A) True; B) False
\\ \textit{Note:} Energy security is interconnected with economic stability, as energy resources are crucial for economic growth and development, and with environmental security, as energy production and consumption significantly impact ecological systems and climate change.
\item \textbf{Answer:} False
\\ \textit{Options:} A) True; B) False
\\ \textit{Note:} The statement is false because despite the United States' technological superiority, the war was prolonged and ultimately unsuccessful due to factors such as guerrilla warfare tactics employed by North Vietnamese forces, a lack of popular support for the war at home, and significant political and social complexities in Vietnam.
\item \textbf{Answer:} False
\\ \textit{Options:} A) True; B) False
\\ \textit{Note:} The serious discussion of cyber-security issues actually began in the 1970 s and 1980 s, predating the widespread popularity of the internet, as concerns around computer security and privacy emerged with the development of early networking technologies.
\end{enumerate}

\section*{Long Form Response}
\begin{enumerate}[label=\arabic*.]
\setcounter{enumi}{10}
\item \textbf{Answer:} The future of security studies faces several key challenges, including the rise of non- state actors, the impact of climate change, and the increasing complexity of cyber threats. These challenges are reshaping global security dynamics by necessitating a more nuanced understanding of security that goes beyond traditional state-centric frameworks. The methodologies employed in the field are also evolving, as scholars are now incorporating interdisciplinary approaches and qualitative analyses to address these multifaceted issues. As security threats become more interlinked and diffuse, the need for adaptive and comprehensive strategies becomes paramount, highlighting the imperative for security studies to evolve in tandem with these emerging challenges. Ultimately, the interplay of these factors will significantly influence both the theoretical foundations and practical applications of security studies in the coming years.
\\ \textit{Note:} correct because it identifies critical emerging challenges in security studies, such as non-state actors, climate change, and cyber threats, and explains how these issues require a shift from traditional state-centric approaches to more interdisciplinary and qualitative methodologies, thereby impacting global security dynamics.
\item \textbf{Answer:} Medicalization is the process by which social issues are increasingly framed as medical problems, leading to a shift in how society understands and responds to these issues. This phenomenon can be observed in various contexts, such as the framing of mental health, addiction, and even social behaviors as medical conditions requiring treatment rather than moral or societal intervention. The implications of medicalization for political and social discourse are significant, as it can lead to a depoliticization of social issues, shifting responsibility from social structures and policies to individual health and behavior. Consequently, this may result in reduced accountability for systemic inequalities and a reliance on medical solutions, overshadowing the need for broader social reform. As a result, the medicalization of social problems complicates public discourse, often limiting the scope of discussion to medical expertise while neglecting the socio-political factors that contribute to these issues.
\\ \textit{Note:} The answer correctly identifies medicalization as the process of framing social issues as medical problems, illustrating its influence on societal understanding and responses, particularly in areas like mental health and addiction, while highlighting the significant implications for political and social discourse, including a potential shift away from moral and societal interventions.
\end{enumerate}

\vfill
\noindent\textit{End of Answer Key}
\end{document}